\section*{第 \ref{ch:b3:relativistic-kinematics} 章——相对论运动学}

\begin{solution}{exo:b3:relativistic-kinematics:1}
$\gamma = 1.005$、$1.155$、$2.294$、$7.09$、$22.4$。国际空间站：$\beta =
\num{2.57e-5}$，$\gamma - 1 \approx \beta^2/2 = \num{3.3e-10}$；在
六个月（$\num{1.6e7}{}\,\unit{s}$）内，宇航员的时钟
\emph{变慢}约 $\approx\qty{5}{ms}$（仅速度效应；
在国际空间站的低轨道高度上，引力效应会减弱，但不会
使这一结论反转）。
\end{solution}

\begin{solution}{exo:b3:relativistic-kinematics:2}
(a) $\Delta\tau = 2d/c$。(b) 每半个周期内，光子走过
斜边：$(c\Delta t/2)^2 = (v\Delta t/2)^2 + d^2$，其中 $d =
c\Delta\tau/2$；解得 $\Delta t = \Delta\tau/\sqrt{1 - v^2/c^2}$。(c)
若 $d$ 随运动而改变，则毕达哥拉斯定理的那一步就不成立。(d)
设两把相同的尺子相向而过，每把尺子的尖端都带着一支画笔，
笔尖指向对方。若运动使横向\omterm{prop:b3:relativistic-kinematics:contraction}{长度收缩}，
则每个参考系都会预言，自己的尺子会在对方尺尖\emph{之外}画出一道痕迹——
两个相互矛盾的事实，针对的是同一相遇\omterm{def:b3:relativistic-kinematics:event}{事件}中的同一组笔迹——
而同一\omterm{def:b3:relativistic-kinematics:event}{事件}中不允许出现矛盾：
横向长度不可能改变。
\end{solution}

\begin{solution}{exo:b3:relativistic-kinematics:3}
$\gamma = 22.4$。(a) $\gamma\tau = \qty{49}{\micro s}$。(b) $v\gamma
\tau = \qty{14.8}{km}$。(c) $15/22.4 = \qty{670}{m}$。(d) 实验室中的飞行时间
为 $\qty{50}{\micro s}$：若不考虑相对论，
$\eu^{-50/2.2} \approx \num{e-10}$，几乎为零；若考虑相对论，
\omterm{prop:b3:relativistic-kinematics:dilation}{固有时}为 $50/22.4 = \qty{2.2}{\micro s}$，比例为 $\eu^{-1}
\approx 0.37$。
\end{solution}

\begin{solution}{exo:b3:relativistic-kinematics:4}
(a) $c\Delta t = \qty{600}{m}$，$\Delta x = \qty{300}{m}$：$\Delta s^2
= (600^2 - 300^2)\,\unit{m^2} > 0$，类时。(b) 可以：
以 $\Delta x/\Delta t = c/2$ 的速率传递的信号即可。(c) 该参考系的速率恰为 $v =
0.5c$；\omterm{prop:b3:relativistic-kinematics:dilation}{固有时}为 $\Delta s/c = \sqrt{600^2 - 300^2}/c = 520/c =
\qty{1.73}{\micro s}$。(d) 此时 $c\Delta t = 300 < \Delta x = 600$：
类空；不存在因果联系；没有任何参考系能使两\omterm{def:b3:relativistic-kinematics:event}{事件}位于同一地点，
但速率为 $v = c^2\Delta t/\Delta x = 0.5c$ 的参考系可使它们
同时发生。
\end{solution}

\begin{solution}{exo:b3:relativistic-kinematics:5}
(a) $\beta = \num{1.29e-5}$：$\gamma - 1 = \num{8.3e-11}$。(b)
$86400\,\unit{s} \times \num{8.3e-11} = \qty{7.2}{\micro s}$ 每天。
(c) $c \times \qty{7.2}{\micro s} = \qty{2.2}{km}$——导航
一天之内便会失效。(d) 净效应 $+38\,\unit{\micro s}$ 表明，
引力蓝移（$+45.7$）的作用超过了运动减慢（
$-7.2$）的作用：在 \qty{20200}{km} 处，位于地球势场中较高的位置，
使时钟走快的作用超过了轨道速度使其变慢的作用。
\end{solution}

\begin{solution}{exo:b3:relativistic-kinematics:6}
(a) $1.6c/1.64 = 0.976c$。(b) $1.8c/1.81 = 0.994c$。(c) $c - u =
\dfrac{(c - u')(c - v)}{c\,(1 + u'v/c^2)}$：两个因子均为正，故
$u < c$。(d) 光斑是一个运动的\emph{图样}，而非实体：没有
任何物质、能量或信息从月面上的一点传到
下一点——每个光子都以 $c$ 的速率射向月球。
\end{solution}

\begin{solution}{exo:b3:relativistic-kinematics:7}
(a) 波长变短：说明正在接近。$f/f_0 = 589/575 = 1.024$；
$(1+\beta)/(1-\beta) = 1.049$：$\beta = 0.024$，约为
\qty{7200}{km/s}。(b) 比值为 $2$：$(1+\beta)/(1-\beta) = 4$，$\beta =
0.6$。(c) 将平方根展开，$\Delta\lambda/\lambda \approx
\beta$；在 \qty{600}{nm} 处，$\qty{1}{\angstrom} = \qty{0.1}{nm}$，
$\beta = \num{1.7e-4}$，即每埃约 \qty{50}{km/s}——天文学家的
条件反射。(d) 横向多普勒效应：纯粹的\omterm{prop:b3:relativistic-kinematics:dilation}{时间延缓}，
量级为 $\beta^2$——只有达到原子级精度
才能测量。
\end{solution}

\begin{solution}{exo:b3:relativistic-kinematics:8}
(a) $20/2 = \qty{10}{m}$：恰好放得下，且两扇门可以同时关闭（谷仓参考系）。
(b) 在跑步者的参考系中，“前门关闭”与“后门打开”这两个
\omterm{def:b3:relativistic-kinematics:event}{事件}\emph{不}同时：出口在入口关闭之前就已打开，
因此长 \qty{20}{m} 的杆子穿过长 \qty{5}{m} 的谷仓时，
始终不会被完全包围。
(c) 谷仓参考系：$\Delta t = 0$，$\Delta x =
\qty{10}{m}$。跑步者参考系：$|\Delta t'| = \gamma v\Delta x/c^2 =
2 \times 0.866c \times 10/c^2 = \qty{58}{ns}$。(d) “杆子被包围”这一说法
隐含地断定两个分离的\omterm{def:b3:relativistic-kinematics:event}{事件}（两扇门都关闭）是同时的，
这是一个依赖参考系的陈述，而非事实。
\end{solution}

\begin{solution}{exo:b3:relativistic-kinematics:9}
$\gamma = 1.25$，对应速率 $0.6c$。(a) 运动中的长度 $100/1.25 = \qty{80}{m}$，
即在 \qty{1.8e8}{m/s} 下为 \qty{444}{ns}。(b) $240/\num{1.8e8} =
\qty{1.33}{\micro s}$。(c) $|\Delta t'| = \gamma v\Delta x/c^2 =
\qty{200}{ns}$；由 $t' = \gamma(t - vx/c^2)$ 可知，位于较大
$x$ 处的夹子——即前方的那个——对火箭而言\emph{先}触发。(d)
空间站参考系：$\Delta s^2 = 0 - 80^2 = -6400\,\unit{m^2}$。火箭参考系：$\Delta
x' = \gamma\,\Delta x = \qty{100}{m}$，$c\Delta t' = \qty{60}{m}$：
$60^2 - 100^2 = -6400\,\unit{m^2}$——不变量。
\end{solution}

\begin{solution}{exo:b3:relativistic-kinematics:10}
(a) 特拉：\qty{20}{yr}；斯特拉：$20/\gamma = \qty{12}{yr}$。(b) 距离收缩为
$8/\gamma = \qty{4.8}{ly}$，她走完该距离所用时间为 $4.8/0.8
= \qty{6}{yr}$（按她自己的时间计），与往返全程的 \qty{12}{yr} 相符。
(c) 与转折\omterm{def:b3:relativistic-kinematics:event}{事件} $(t =
\qty{10}{yr},\ x = \qty{8}{ly})$ 同时的地球时钟：去程参考系给出
$t - vx/c^2 = 10 - 6.4 = \qty{3.6}{yr}$；回程参考系给出 $10 +
6.4 = \qty{16.4}{yr}$。她所认为的“地球上此刻”跳跃了 \qty{12.8}{yr}
（在转折处）；她的账目 $3.6 + 12.8 + 3.6 = \qty{20}{yr}$ 与
特拉完全吻合。(d) 斯特拉占据\emph{两个}\omterm{thm:b3:relativistic-kinematics:postulates}{惯性参考系}，它们之间由
她能感受到的加速度相连；特拉只占据一个。连接出发与重逢两\omterm{def:b3:relativistic-kinematics:event}{事件}的
直线型世界线所携带的\omterm{prop:b3:relativistic-kinematics:dilation}{固有时}最长，
只有斯特拉的世界线是弯曲的。
\end{solution}

\begin{solution}{exo:b3:relativistic-kinematics:11}
(a) 若映射是非线性的，则匀速运动在另一参考系中将不再呈匀速，
从而无法区分均匀空间与均匀时间中的各点。
(b) 将一个关系式代入另一个：$t' = \gamma t +
(1 - \gamma^2)x/\gamma v$。(c) 令 $x = ct$，$x' = ct'$：
$\gamma(c - v)t = c\gamma t + c(1 - \gamma^2)ct/\gamma v$，由此解得
$\gamma^2 = 1/(1 - v^2/c^2)$。(d) 相对性给出两个方向上相同的
$\gamma$（不存在优选参考系）；光速公设将剩下的自由函数
确定为
$\gamma(v)$。
\end{solution}

\begin{solution}{exo:b3:relativistic-kinematics:12}
(a) 令 $\cosh\varphi = \gamma$，$\sinh\varphi = \gamma\beta$，
变换恰为所述的双曲旋转，且
$\cosh^2 - \sinh^2 = 1$ 保持 $c^2t^2 - x^2$ 不变。(b)
$\tanh(\varphi_1 + \varphi_2) = (\tanh\varphi_1 + \tanh\varphi_2)/(1 +
\tanh\varphi_1\tanh\varphi_2)$：恰为\omterm{prop:b3:relativistic-kinematics:addition}{速度合成}律——
速度不相加，快度才相加。(c) 由 $\varphi = g\tau/c$，
得 $\beta = \tanh(g\tau/c)$；由 $\operatorname{artanh}(0.99) = 2.65$ 得
$\tau = 2.65c/g \approx \qty{8.1e7}{s} \approx 2.6$ 年（飞船时间）。(d) $\varphi$ 取遍全体实数，
而 $\tanh\varphi$ 趋于饱和值
$1$：可以无限加速下去，快度线性增长，
而速率只是无限趋近于 $c$。
\end{solution}

\begin{solution}{pb:b3:relativistic-kinematics:1}
\textbf{1.} 所有 μ 子完全相同，并以固定的统计速率衰变：
一群体的存活比例度量经过的\omterm{prop:b3:relativistic-kinematics:dilation}{固有时}，
其准确程度不亚于钟表的指针。
\textbf{2.} $c\tau = \qty{3.00e8}{} \times \qty{2.2e-6}{} =
\qty{660}{m}$.
\textbf{3.} $t = \qty{15}{km}/c = \qty{50}{\micro s}$：存活比例
$\eu^{-50/2.2} = \eu^{-22.7} \approx \num{1.4e-10}$.
\textbf{4.} 本“不可能”穿行超过一千米的粒子，
却穿过了其中十五个，并大量到达。
\textbf{5.} 一只手掌 $\sim\qty{100}{cm^2}$：每秒约有几个 μ 子；
人体水平截面 $\sim\num{e3}\,\unit{cm^2}$：
两次心跳之间约有十五个——相对论的“雨”
无时无刻不穿过每一个人。
\textbf{6.} $\gamma = 1/\sqrt{1 - 0.995^2} = 10.0$.
\textbf{7.} 比较\emph{同一}被选群体在已知高度差上的
两次计数，即可消去
关于 μ 子产生于何处、产生多少的一切假设。
\textbf{8.} $t = 1907/(0.995 \times \num{3e8}) = \qty{6.39}{\micro s}$.
\textbf{9.} $\eu^{-6.39/2.2} = 0.055$：约每小时 $31$ 个。
\textbf{10.} $t/\gamma = \qty{0.64}{\micro s}$.
\textbf{11.} $\eu^{-0.64/2.2} = 0.75$：约每小时 $420$ 个。
\textbf{12.} 实测 $408 \pm 9$：相对论的 $\approx 420$ 在
μ 子于探测器吸收体中轻微减速所允许的范围内与之相符；
经典物理的 $31$ 则偏差约十三倍。山
作出了裁决。
\textbf{13.} $408/563 = 0.725 = \eu^{-t_{\text{固有}}/\tau}$：
故 $t_{\text{固有}} = \qty{0.71}{\micro s}$，$\gamma_{\exp} =
6.39/0.71 = 9.0$——与弗里施和史密斯的 $8.8 \pm 0.8$ 十分吻合。
\textbf{14.} $1907/10 = \qty{191}{m}$.
\textbf{15.} $191/(0.995c) = \qty{0.64}{\micro s}$：同样的 $25\%$
衰变——这是 μ 子对同一计数的描述。
\textbf{16.} 一次计数是一组局域的符合\omterm{def:b3:relativistic-kinematics:event}{事件}；
\omterm{def:b3:relativistic-kinematics:event}{事件}及其计数在参考系变换下不变——
参考系对时间与长度可以有不同意见，但对计数器读到的内容永不会有分歧。
\textbf{17.} $\Delta s^2 = (c \times \qty{6.39}{\micro s})^2 -
(\qty{1907}{m})^2 = (1917^2 - 1907^2)\,\unit{m^2} = (\qty{196}{m})^2$:
$\Delta s/c = \qty{0.65}{\micro s}$——即\omterm{prop:b3:relativistic-kinematics:dilation}{固有时}，
无需离开山上的参考系即可算出。
\textbf{18.} 若衰变率只取决于海拔，则对所有速率都相同；
而实际上，存活比例随 $\gamma$ 变化——速率较慢的被选群体\omterm{prop:b3:relativistic-kinematics:dilation}{时间延缓}较小，
恰与 $\gamma(v)$ 的预言一致。
\textbf{19.} 寿命 $\gamma\tau$ 精确到千分之一（CERN 的 μ 子储存环）；
环形装置使 μ 子成为一个持续加速的旅行者——
“旅行的孪生子”仍保持年轻，
即使旅程是一场没有尽头的掉头。
\textbf{20.} $\beta = \num{8.3e-7}$: $\Delta t = \tfrac12\beta^2
\times \qty{144000}{s} = \qty{50}{ns}$.
\textbf{21.} 飞机的速率与旋转地球的速率相加或相减，
而高度（引力）效应在两个方向上相同：
向东与向西的飞行将两项贡献清楚地分开
（1971 年的结果与两者均相符）。
\textbf{22.} 由\cref{exo:b3:relativistic-kinematics:5}可知：
每天 \qty{7.2}{\micro s}。
\textbf{23.} $7 \times 7.2\,\unit{\micro s} \times c \approx
\qty{15}{km}.$
\textbf{24.} 位置会在数小时内漂移数百米，
在一天之内漂移数千米——导航
将在第一个早晨便明显失效。
\textbf{25.} 1963 年，一座山、两台计数器与 \qty{2.2}{\micro s} 的钟
测得 $\gamma_{\exp} = 9.0$，而预言值为 $10$（计入减速后为 $8.8 \pm 0.8$）：
\omterm{prop:b3:relativistic-kinematics:dilation}{时间延缓}在 $10\%$ 的范围内得到证实——
而同样的物理在你头顶上方的每一颗导航卫星中，
每一秒都被悄然校正。
\end{solution}
